\documentclass[conference,onecolumn]{IEEEtran}

\usepackage[T1]{fontenc}
\usepackage[utf8]{inputenc}
\usepackage[spanish,es-tabla]{babel}
\usepackage{lmodern}
\usepackage{microtype}
\usepackage{csquotes}
\usepackage{amsmath,amssymb}
\usepackage{booktabs}
\usepackage{float}
\usepackage{tabularx}
\usepackage{array}
\usepackage{url}
\usepackage[hidelinks]{hyperref}
\usepackage[backend=biber,style=ieee,sorting=none]{biblatex}
\addbibresource{referencias.bib}
\AtBeginDocument{\spanishplainpercent}

\newcolumntype{Y}{>{\raggedright\arraybackslash}X}
\newcommand{\si}{\ensuremath{\checkmark}}
\newcommand{\no}{\ensuremath{\times}}
\newcommand{\req}[1]{\par\noindent\textbf{#1.}\hspace{0.5em}}
\renewcommand{\IEEEkeywordsname}{Palabras clave}

\title{Aplicación de ISO/IEC/IEEE 29148:2018 a la Especificación de Requisitos de un Sistema Bibliotecario en Línea para la UTEQ}
\author{\IEEEauthorblockN{Allan, Danela, Winston, Gilces}\IEEEauthorblockA{Carrera de Ingeniería de Software\\
Facultad de Ciencias de la Computación\\
Universidad Técnica Estatal de Quevedo, Ecuador\\
Ingeniería de Requisitos -- ISR-401}}

\begin{document}
\maketitle

\begin{abstract}
Avance académico acumulativo de la práctica experimental del sistema bibliotecario en línea de la UTEQ. Contenido incorporado: 5 stakeholders, 4 RF. Los requisitos y criterios aún requieren validación con los interesados.
\end{abstract}

\begin{IEEEkeywords}
Ingeniería de requisitos, ISO/IEC/IEEE 29148, UTEQ, sistema bibliotecario.
\end{IEEEkeywords}

\section{Tarea 1: Identificación de stakeholders}
Los \emph{stakeholders} son personas, organizaciones o entidades que participan en el sistema, influyen en sus decisiones o se ven afectadas por su funcionamiento; su identificación y consulta constituyen parte del proceso de Ingeniería de Requisitos \cite{ieee29148,pohl2025}. En este primer avance se identifican cinco interesados internos, con niveles de influencia estimados para esta práctica (Tabla~\ref{tab:stakeholders}).

\begin{table}[H]
\caption{Stakeholders del sistema bibliotecario de la UTEQ (identificación preliminar)}
\label{tab:stakeholders}
\centering\footnotesize
\begin{tabularx}{\textwidth}{@{}p{0.22\textwidth}p{0.14\textwidth}Yp{0.09\textwidth}@{}}
\toprule
\textbf{Stakeholder} & \textbf{Rol} & \textbf{Interés principal en el sistema} & \textbf{Influencia}\\
\midrule
Estudiantes de las 17 carreras & Interno, usuario final & Consultar disponibilidad, reservar libros y controlar sus préstamos. & Medio\\
Docentes universitarios & Interno, usuario final & Acceder a bibliografía para actividades académicas y de investigación. & Medio\\
Bibliotecarios & Interno, operador & Registrar préstamos, devoluciones, renovaciones y administrar ejemplares. & Alto\\
Responsable de biblioteca & Interno, administrador & Supervisar el catálogo, controlar la circulación y obtener estadísticas. & Alto\\
Administradores de TI & Interno, técnico & Garantizar seguridad, autenticación institucional, integración y funcionamiento del sistema. & Alto\\
\bottomrule
\end{tabularx}
\end{table}

\section{Tarea 2: Diez requisitos funcionales}
Los requisitos funcionales describen capacidades y comportamientos observables del sistema. La redacción busca satisfacer las propiedades de un requisito individual establecidas por ISO/IEC/IEEE 29148:2018 \cite{ieee29148}; todos los requisitos siguen la plantilla exigida en el enunciado del curso \cite{uteqsemana1}:

\begin{quote}
\textbf{RF-NN:} El sistema deberá [acción] [objeto] cuando [condición], con el fin de [propósito].
\end{quote}

\subsection{Requisitos propuestos}
\req{RF-01: Registro y acceso institucional}El sistema deberá habilitar el perfil bibliotecario del usuario, creado automáticamente si todavía no existe, cuando este inicie sesión con su identidad institucional validada, con el fin de identificarlo para utilizar los servicios de préstamo.
\req{RF-02: Búsqueda de libros}El sistema deberá mostrar los libros coincidentes con el título, autor o ISBN consultado, incluyendo su disponibilidad, cuando el usuario ingrese un criterio de búsqueda, con el fin de localizar material bibliográfico.
\req{RF-03: Reserva de libros}El sistema deberá registrar la reserva de un título sin ejemplares disponibles cuando un usuario autenticado la solicite sin tener otra reserva activa del mismo título, con el fin de asignarle un turno de espera.
\req{RF-04: Registro de préstamos}El sistema deberá registrar el préstamo de un ejemplar disponible con vencimiento a siete días calendario cuando el bibliotecario confirme su entrega a un usuario con menos de tres préstamos activos y sin restricciones vigentes, con el fin de controlar la circulación del ejemplar.

\printbibliography[title={Referencias}]
\end{document}
